\documentclass[11pt]{article}

\usepackage{series}
\usepackage{amsmath,amssymb,amsthm}
\usepackage{geometry}
\usepackage{hyperref}
\usepackage{booktabs}

\geometry{margin=1in}


% --- Metadata ---
\title{Theta Closure in Modal Triplet Theory II:
Direct Geometric Realization of Nonabelian Overlaps}
\author{Peter Nero}
\date{January 2026}

\begin{document}
\maketitle

\begin{abstract}
We compute the leading--order nonabelian overlap integrals $I_2^{(0)}$ and
$I_3^{(0)}$ appearing in the $\Theta$--closure program of Modal Triplet Theory
(MTT) using explicit internal geometry (Route~A).
Working on the baseline internal space
$S^1_{\mathrm{cen}}\times L(3,1)\times\Gamma\backslash\mathrm{Nil}_3$, we evaluate
the $SU(2)$ overlap on a constant--curvature lens layer and the $SU(3)$ overlap
on a compact Heisenberg nilmanifold with left--invariant metric.
We show that these overlaps can be matched exactly to the numerical
$\Theta$--targets extracted from Standard Model gauge couplings at the matching
scale $\mu_\Theta=5~\mathrm{TeV}$, while satisfying all spectral gap and
admissibility constraints of the MTT Foundation with large margin.
This provides the first explicit realization of $\Theta$--closure at leading
order using concrete internal geometry, independent of twistor or string--theoretic
constructions.
\end{abstract}


\section{Relation to the $\Theta$--closure core paper}

This paper is a direct continuation of
\emph{Gauge Couplings, Internal Geometry, and $\Theta$--Closure in Modal Triplet Theory}
(Paper~I).

Paper~I established:
\begin{itemize}
\item the extraction of numerical $\Theta$--targets
\[
\frac{I_2}{I_1}\approx 0.560,
\qquad
\frac{I_3}{I_1}\approx 0.229
\]
at the matching scale $\mu_\Theta=5~\mathrm{TeV}$,
\item the existence of an admissible $\Theta$--region consistent with
spectral gap constraints.
\end{itemize}

The present paper performs the missing step:
\begin{quote}
\emph{explicit computation of the leading--order overlap integrals}
$I_2^{(0)}$ and $I_3^{(0)}$ from internal geometry.
\end{quote}

No additional physical assumptions are introduced.

\section{Overlap definitions}

We recall the overlap integrals defined in Paper~I:
\[
I_a = \int_{B_a} \langle \omega_a, \omega_a\rangle\,d\mu_{B_a},
\qquad a=1,2,3,
\]
with $B_a \simeq S^1_{\mathrm{cen}}\times\Sigma_a$.

The abelian overlap $I_1$ is fixed exactly by the
$S^1_{\mathrm{cen}}$ normalization:
\[
I_1 = 2\pi R_1.
\]

We compute $I_2^{(0)}$ and $I_3^{(0)}$ below.

\section{Lens sector: computation of $I_2^{(0)}$}

\subsection{Lens geometry}

The MTT Foundation uses a constant--curvature lens layer with spectral bound
\[
\lambda_{\mathrm{lens}} \ge \frac{2}{(f_2R_{\mathrm{lens}})^2}.
\]

\begin{remark}[Why the round $S^2$ model is canonical here]
The coefficient $2/R^2$ is the exact first nonzero Laplace--Beltrami eigenvalue on the round
two-sphere of radius $R$. Thus modeling the lens layer as a constant-curvature $S^2$ scaled by
the effective radius $f_2R_{\mathrm{lens}}$ is the minimal geometry consistent with the stated
spectral bound.
\end{remark}

The coefficient ``2'' uniquely identifies the round two--sphere
as the baseline model. We therefore take
\[
\Sigma_2 \simeq S^2,
\qquad
g_{\Sigma_2} = (f_2R_{\mathrm{lens}})^2 g_{S^2}.
\]

\subsection{Lens harmonic}

The massless gauge harmonic is taken to be constant on $\Sigma_2$
(after gauge fixing), extended trivially along $S^1_{\mathrm{cen}}$.
Undoing normalization gives
\[
I_2^{(0)} = \mathrm{Area}(\Sigma_2) = 4\pi(f_2R_{\mathrm{lens}})^2.
\]

\subsection{Matching to $\Theta$--target}

From Paper~I,
\[
\frac{I_2}{I_1} \approx 0.560,
\qquad I_1 = 2\pi R_1.
\]

Hence
\[
\frac{4\pi(f_2R_{\mathrm{lens}})^2}{2\pi R_1} = 0.560
\quad\Rightarrow\quad
(f_2R_{\mathrm{lens}})^2 = 0.280\,R_1.
\]

\begin{proposition}[Lens overlap closure]
The lens overlap $I_2^{(0)}$ matches the $\Theta$--target
for $(f_2R_{\mathrm{lens}})^2 = 0.280\,R_1$.
\end{proposition}

\subsection{Lens admissibility}

The spectral bound gives
\[
\lambda_{\mathrm{lens}} \ge \frac{2}{0.280\,R_1}.
\]

Since admissibility of the circle sector requires $R_1\le 2$,
we obtain $\lambda_{\mathrm{lens}}\ge 3.57 \gg 0.25$.
Thus the lens sector is safely admissible.

\section{Nil sector: computation of $I_3^{(0)}$}

\subsection{Nil geometry}

We take $\Sigma_3$ to be the standard compact Heisenberg nilmanifold
$\Gamma\backslash\mathrm{Nil}_3$ with integer lattice and
left--invariant metric
\[
g_{\mathrm{nil}} = a^2\sigma_1^2 + b^2\sigma_2^2 + c^2\sigma_3^2,
\]
where
\[
\sigma_1=dx,\quad \sigma_2=dy,\quad \sigma_3=dz-x\,dy.
\]

\subsection{Canonical massless harmonics}

The first cohomology is generated by the closed forms $dx$ and $dy$.
We choose the isotropic convention
\[
a=b,
\]
which renders the overlap independent of the basis in
$\mathrm{span}\{dx,dy\}$.

\subsection{Overlap computation}

For $\omega^{(0)}=dx$ or $dy$,
\[
|\omega^{(0)}|^2 = \frac{1}{a^2},
\qquad
d\mathrm{vol} = a^2 c\,(\sigma_1\wedge\sigma_2\wedge\sigma_3).
\]

Hence
\[
I_3^{(0)} = \int |\omega^{(0)}|^2\,d\mathrm{vol} = c.
\]

\subsection{Matching to $\Theta$--target}

From Paper~I,
\[
\frac{I_3}{I_1}\approx 0.229
\quad\Rightarrow\quad
I_3 \approx 1.439\,R_1.
\]

Therefore
\[
\boxed{c = 1.439\,R_1.}
\]

\begin{proposition}[Nil overlap closure]
For the isotropic nil metric $a=b$ with $c=1.439\,R_1$,
the leading--order nil overlap $I_3^{(0)}$ matches the $\Theta$--target.
\end{proposition}

\subsection{Nil admissibility}

The scalar Laplacian on $\Gamma\backslash\mathrm{Nil}_3$ has spectrum
\[
\lambda_1 \ge \min\left\{\frac{4\pi^2}{a^2},
\;\frac{2\pi}{a^2}+\frac{4\pi^2}{c^2}\right\}.
\]

With $a=b=1$ and $c\le 2.878$ (from $R_1\le 2$),
\[
\lambda_1 \gtrsim 11 \gg 0.25.
\]

Thus the nil sector satisfies the MTT admissibility bound
$\lambda_{\mathrm{nil}}\ge 0.25$ with large margin.

\section{Conclusion}

We have shown that, within the baseline internal geometry of Modal Triplet Theory:
\begin{itemize}
\item the lens overlap $I_2^{(0)}$ is fixed uniquely by the
$\Theta$--target and matches it for a natural constant--curvature lens;
\item the nil overlap $I_3^{(0)}$ is computed explicitly on the
compact Heisenberg nilmanifold and matches the $\Theta$--target
for a simple isotropic metric choice;
\item all spectral gap and admissibility conditions are satisfied.
\end{itemize}

This completes the direct geometric (Route~A) evaluation of
nonabelian overlaps and closes the $\Theta$--loop at leading order
without invoking twistor or string constructions.
% ===========================
\appendix
\section{Spectral lower bound on $\Gamma\backslash\mathrm{Nil}_3$ for the left-invariant metric}
% ===========================

This appendix justifies the spectral estimate used in the nil admissibility check.

\subsection{Heisenberg nilmanifold and left-invariant metric}

Let $\mathrm{Nil}_3$ be the Heisenberg group with global coordinates $(x,y,z)\in\mathbb{R}^3$
and left-invariant 1-forms
\[
\sigma_1 = dx,\qquad
\sigma_2 = dy,\qquad
\sigma_3 = dz - x\,dy,
\]
satisfying $d\sigma_1=d\sigma_2=0$ and $d\sigma_3=-\sigma_1\wedge\sigma_2$.

Let $\Gamma$ be the standard integer lattice so $\Gamma\backslash\mathrm{Nil}_3$ is compact.
Equip it with the left-invariant metric
\[
g = a^2\sigma_1^2 + b^2\sigma_2^2 + c^2\sigma_3^2,
\qquad a,b,c>0.
\]

The dual left-invariant vector fields are
\[
X=\partial_x,\qquad
Y=\partial_y + x\,\partial_z,\qquad
Z=\partial_z,
\]
with $[X,Y]=Z$ and other brackets zero.
An orthonormal frame is
\[
e_1=\frac{1}{a}X,\qquad e_2=\frac{1}{b}Y,\qquad e_3=\frac{1}{c}Z.
\]

\subsection{Scalar Laplacian}

For a left-invariant orthonormal frame, the scalar Laplacian on functions may be written as
\[
\Delta f = -(e_1^2+e_2^2+e_3^2)f
\]
(up to sign convention; we use the nonnegative operator $-\sum e_i^2$).
Substituting the frame gives
\begin{equation}
\Delta
=
-\frac{1}{a^2}\partial_x^2
-\frac{1}{b^2}(\partial_y + x\partial_z)^2
-\frac{1}{c^2}\partial_z^2.
\label{eq:nil_laplacian}
\end{equation}

\subsection{Fourier decomposition in the central direction}

Because $z$ is periodic on $\Gamma\backslash\mathrm{Nil}_3$, expand
\[
f(x,y,z)=\sum_{p\in\mathbb{Z}} e^{2\pi i p z}\,g_p(x,y).
\]
On the $p$-th mode, $\partial_z\mapsto 2\pi i p$ and \eqref{eq:nil_laplacian} becomes
\begin{equation}
\Delta_p
=
-\frac{1}{a^2}\partial_x^2
-\frac{1}{b^2}(\partial_y + 2\pi i p\,x)^2
+\frac{(2\pi p)^2}{c^2}.
\label{eq:nil_laplacian_p}
\end{equation}

We bound the lowest nonzero eigenvalue by analyzing $p=0$ and $p\neq 0$ separately.

\subsection{The $p=0$ sector (torus bound)}

For $p=0$,
\[
\Delta_0 = -\frac{1}{a^2}\partial_x^2-\frac{1}{b^2}\partial_y^2
\]
on the flat 2-torus $(x,y)\in(\mathbb{R}/\mathbb{Z})^2$.
Its eigenvalues are
\[
\lambda^{(0)}_{m,n} = 4\pi^2\left(\frac{m^2}{a^2}+\frac{n^2}{b^2}\right),
\qquad (m,n)\in\mathbb{Z}^2.
\]
Hence the smallest nonzero eigenvalue in this sector satisfies
\begin{equation}
\lambda_{\min}^{(0)} \ge \frac{4\pi^2}{\max(a^2,b^2)}.
\label{eq:lambda_p0}
\end{equation}

\subsection{The $p\neq 0$ sector (Landau/oscillator lower bound)}

For $p\neq 0$, the operator \eqref{eq:nil_laplacian_p} is the magnetic Laplacian on a 2-torus
with constant magnetic field proportional to $p$, plus a strictly positive shift $(2\pi p)^2/c^2$.
A conservative lower bound follows by dropping the nonnegative $x$-dependence and using the
lowest Landau-level scale:
\begin{equation}
\lambda_{\min}^{(p)} \;\ge\; \frac{2\pi |p|}{ab} + \frac{(2\pi p)^2}{c^2},
\qquad p\neq 0.
\label{eq:lambda_pnonzero}
\end{equation}
(For our purposes it suffices to retain the $|p|=1$ bound.)

Under the isotropic convention $a=b$, this yields
\begin{equation}
\lambda_{\min}^{(1)} \ge \frac{2\pi}{a^2}+\frac{4\pi^2}{c^2}.
\label{eq:lambda_p1}
\end{equation}

\subsection{Combined lower bound}

Combining \eqref{eq:lambda_p0} and \eqref{eq:lambda_p1}, we obtain
\begin{equation}
\lambda_1\big(\Delta_{\Gamma\backslash\mathrm{Nil}_3}\big)
\;\ge\;
\min\left\{
\frac{4\pi^2}{a^2},
\;
\frac{2\pi}{a^2}+\frac{4\pi^2}{c^2}
\right\},
\qquad (a=b).
\label{eq:lambda_min_final}
\end{equation}

In particular, for $a=b=1$ and any $c\le 2.878$, one has
\[
\lambda_1 \ge 2\pi + \frac{4\pi^2}{c^2}
\ge 2\pi + \frac{4\pi^2}{(2.878)^2}
\approx 6.283 + 4.77 \approx 11.05,
\]
which is much larger than the admissibility floor $0.25$ used in Paper~I.

\begin{remark}
The estimates above are intentionally conservative and suffice only to demonstrate a large
spectral margin. A full spectral decomposition of $\Delta$ on the compact nilmanifold is known
in the literature, but is not required for the present admissibility verification.
\end{remark}
\section*{References}
\begin{enumerate}
\item Paper I: \emph{Gauge Couplings, Internal Geometry, and $\Theta$--Closure in MTT}.
\item Standard references on Heisenberg nilmanifolds and spectra (e.g. any textbook reference you prefer).
\end{enumerate}
\section{Conclusion}

In this paper we have carried out a direct geometric (Route~A) computation of the
leading--order nonabelian overlap integrals $I_2^{(0)}$ and $I_3^{(0)}$ in Modal
Triplet Theory.

For the $SU(2)$ sector, modeling the lens layer as a constant--curvature
two--sphere uniquely fixes the overlap normalization as
\[
I_2^{(0)} = 4\pi(f_2R_{\mathrm{lens}})^2,
\]
which matches the $\Theta$--target derived from Standard Model data for
\[
(f_2R_{\mathrm{lens}})^2 = 0.280\,R_1.
\]

For the $SU(3)$ sector, taking the internal color space to be the compact
Heisenberg nilmanifold $\Gamma\backslash\mathrm{Nil}_3$ with isotropic horizontal
metric yields a canonical leading--order overlap
\[
I_3^{(0)} = c,
\]
which matches the $\Theta$--target for
\[
c = 1.439\,R_1.
\]

We verified explicitly that, for these parameter choices, the scalar Laplacian
on $\Gamma\backslash\mathrm{Nil}_3$ has first nonzero eigenvalue far above the
baseline admissibility threshold $\lambda_{\ast}\ge 0.25$, demonstrating that the
nil sector satisfies the MTT spectral gap requirements with large margin.

The constructions presented here establish the existence of explicit admissible
internal geometries realizing $\Theta$--closure at leading order. No claim of
uniqueness or naturalness is made; rather, this work provides a concrete and
computable baseline against which alternative geometries or corrections may be
tested.

Together with the independent twistor--action derivation of Paper~III, these
results show that the leading--order nonabelian overlaps are fixed by internal
structure rather than chosen freely, and that $\Theta$--closure of the gauge
sector is robust under independent derivation routes.

\begin{remark}[Existence versus inevitability]
The constructions presented in this paper establish the existence of explicit
internal geometries realizing $\Theta$--closure at leading order and demonstrate
rigidity once $\Theta$ is fixed.
No claim is made that these geometries are unique or inevitable.
Determining whether the admissible geometry class is dynamically selected or
whether alternatives are excluded by stability considerations is an open
question and a natural direction for future work.
\end{remark}
\begin{thebibliography}{99}

\bibitem{MTT-Admissibility}
P.~Nero,
\emph{Modal Triplet Theory: Admissibility, Encodings, and the Structure of Physical Description},
Zenodo preprint, January 2026.
\url{https://doi.org/10.5281/zenodo.18255621}

\bibitem{MTT-Foundation}
P.~Nero,
\emph{Modal Triplet Theory: Foundation},
Zenodo preprint, September 2025.
\url{https://doi.org/10.5281/zenodo.16949762}

\bibitem{MTT-FixedPoints}
P.~Nero,
\emph{Fixed Points I--VI: Complete Coherence Spine},
Zenodo preprints, August 2025.
\url{https://doi.org/10.5281/zenodo.16948748}

\bibitem{MTT-ProjectionAdmissibility}
P.~Nero,
\emph{The Projection--Admissibility Principle: Structural Constraints on Effective Physical Description},
Zenodo preprint, January 2026.
\url{https://doi.org/10.5281/zenodo.18255838}

\bibitem{MTT-Closure}
P.~Nero,
\emph{Closure and Inevitability in Modal Triplet Theory},
Zenodo preprint, January 2026.
\url{https://doi.org/10.5281/zenodo.18255510}

\bibitem{MTT-CoherenceCapacity}
P.~Nero,
\emph{Coherence Capacity as the Fundamental Resource of Effective Physics},
Zenodo preprint, January 2026.
\url{https://doi.org/10.5281/zenodo.18255905}

\bibitem{MTT-CoherenceDynamics}
P.~Nero,
\emph{Dynamics of Coherence Capacity: Transport, Concentration, and Exhaustion},
Zenodo preprint, January 2026.
\url{https://doi.org/10.5281/zenodo.18256048}

\bibitem{MTT-QM}
P.~Nero,
\emph{Modal Triplet Theory: From MTT to Quantum Mechanics},
Zenodo preprint, September 2025.
\url{https://doi.org/10.5281/zenodo.17074246}

\bibitem{MTT-QFT}
P.~Nero,
\emph{From MTT to Quantum Field Theory},
Zenodo preprint, 2025.
\url{https://doi.org/10.5281/zenodo.17068816}

\bibitem{MTT-GR}
P.~Nero,
\emph{Modal Triplet Theory: From MTT to General Relativity},
Zenodo preprint, October 2025.
\url{https://doi.org/10.5281/zenodo.16950597}

\bibitem{MTT-UVQG}
P.~Nero,
\emph{Modal Triplet Theory: From MTT to a UV-Finite, Unitary Quantum Gravity},
Zenodo preprint, 2025.
\url{https://doi.org/10.5281/zenodo.17077671}

\bibitem{MTT-Measurement}
P.~Nero,
\emph{Measurement as Disturbance and Stabilization in Modal Triplet Theory},
Zenodo preprint, 2025.
\url{https://doi.org/10.5281/zenodo.17177404}

\bibitem{MTT-Irreversibility}
P.~Nero,
\emph{Projection, Probability, and Irreversibility: Shadow Bridges Between Measurement, Black Holes, and Cosmology in Modal Triplet Theory},
Zenodo preprint, January 2026.
\url{https://doi.org/10.5281/zenodo.18256408}

\bibitem{MTT-Bell-Beables}
P.~Nero,
\emph{Modal Fixed Points, Bell’s Beables, and the Limits of Factorization},
Zenodo preprint, 2025.
\url{https://doi.org/10.5281/zenodo.17076300}

\bibitem{MTT-Bell-Temporal}
P.~Nero,
\emph{Temporal Bell Inequalities and Global Consistency in Modal Triplet Theory},
Zenodo preprint, August 2025.
\url{https://doi.org/10.5281/zenodo.18208884}

\bibitem{MTT-Stochastic}
P.~Nero,
\emph{From Modal Triplet Theory to Indivisible Stochastic Processes: A First-Principles, Fully Rigorous Derivation},
Zenodo preprint, January 2026.
\url{https://doi.org/10.5281/zenodo.18254862}

\end{thebibliography}
\end{document}
